\chapter{Prompts y plantillas}
\label{ap:prompts}
% Fuentes: Codigo/scripts/train_vlm_classify.py (CLASSIFICATION_PROMPT); Codigo/src/evaluation/vlm_scoring.py
% (CLASS_SYNONYMS, PROMPT_TEMPLATES, COT_PROMPT); Codigo/src/models/dermlip.py (DEFAULT_TEMPLATES);
% Codigo/scripts/generate_report.py (REPORT_PROMPT, GROUNDED_TEMPLATES, RAG_TEMPLATE, GROUNDED_RAG_TEMPLATE,
% DIFFERENTIAL_TEMPLATE); Codigo/scripts/dermlip_split_eval.py (frase de recuperación);
% Codigo/src/evaluation/clinical_ner.py (DERMATOLOGY_TERMS, CLASS_TERMS).

Este apéndice reúne los textos que reciben los modelos, en inglés y tal como aparecen en el código. Los campos
entre llaves se completan en cada imagen: \texttt{\{dx\}} con el nombre completo de la clase, como
«melanoma», \texttt{\{candidates\}} con los dos candidatos separados por punto y coma, y
\texttt{\{rag\_context\}} con la frase de los casos recuperados.

\section{Clasificación con modelos visión--lenguaje generativos}

La estrategia \emph{zero-shot} por generación (Sección~\ref{sec:met_vlm_estrategias}) usa una sola instrucción,
que pide responder con el código de la clase:
\begin{minted}[linenos=false]{text}
Classify this dermoscopic image into one of the following categories: MEL, NV,
BCC, AKIEC, BKL, DF, VASC. Respond with only the category label.
\end{minted}

En cambio, la estrategia de verosimilitud promedia cuatro \emph{prompts} (Ecuación~\ref{eq:met_likelihood}). Los
dos primeros son abiertos, porque no nombran las clases, y los dos últimos las enumeran y terminan justo donde
debería ir el nombre de la clase:
\begin{minted}[linenos=false]{text}
O1:
    What is the most likely diagnosis for this dermoscopic image?
O2:
    Identify the type of skin lesion shown in this image.
C1:
    Given this dermoscopic image, the diagnosis among melanoma, melanocytic
    nevus, basal cell carcinoma, actinic keratosis, benign keratosis,
    dermatofibroma, or vascular lesion is
C2:
    Select the correct diagnosis from: melanoma, melanocytic nevus, basal cell
    carcinoma, actinic keratosis, benign keratosis, dermatofibroma, vascular
    lesion. Answer:
\end{minted}
Cada clase se puntúa con tres sinónimos, y su puntaje es el máximo (Tabla~\ref{tab:ap_sinonimos}).

\begin{table}[H]
    \centering
    \caption{Sinónimos con que se puntúa cada clase en la estrategia de verosimilitud.}
    \label{tab:ap_sinonimos}
    \small
    \begin{tabularx}{\textwidth}{@{}>{\raggedright\arraybackslash}p{2.2cm}>{\raggedright\arraybackslash}X@{}}
        \toprule
        \textbf{Clase} & \textbf{Sinónimos} \\
        \midrule
        MEL & melanoma; malignant melanoma; MEL \\
        \filasep
        NV & melanocytic nevus; nevus; mole \\
        \filasep
        BCC & basal cell carcinoma; basal cell; BCC \\
        \filasep
        AKIEC & actinic keratosis; intraepithelial carcinoma; AKIEC \\
        \filasep
        BKL & benign keratosis; seborrheic keratosis; solar lentigo \\
        \filasep
        DF & dermatofibroma; fibrous histiocytoma; DF \\
        \filasep
        VASC & vascular lesion; angioma; VASC \\
        \bottomrule
    \end{tabularx}
\end{table}

Por último, el razonamiento paso a paso fija un formato de tres campos, y la clase se lee de la línea final:
\begin{minted}[linenos=false]{text}
You are a dermatology expert analyzing a dermoscopic image. Think step by step
and use the following format exactly:
Features: [describe asymmetry, borders, colors, diameter, special structures]
Differential: [list the most likely diagnoses with brief reasoning]
Final diagnosis: [one of: MEL, NV, BCC, AKIEC, BKL, DF, VASC]
\end{minted}

\section{Clasificación con DermLIP}

El modo \emph{zero-shot} de DermLIP (Sección~\ref{sec:met_dermlip}) sigue el protocolo de sus
autores~\cite{yan2025derm1miccv}: promedia, para cada clase, los \emph{embeddings} de texto de ocho plantillas, en las
que \texttt{\{c\}} es el nombre de la clase (melanoma, nevus, basal cell carcinoma, actinic keratosis, benign keratosis,
dermatofibroma y vascular):
\begin{minted}[linenos=false]{text}
This is a skin image of {c}
This is a skin image of {c}.
A skin image of {c}.
An image of {c}, a skin condition.
{c}, a skin disorder, is shown in this image.
The skin lesion depicted is {c}.
The skin cancer in this image is {c}.
This image depicts {c}, a type of skin cancer.
\end{minted}

\section{Generación de reportes}

Tanto el ajuste en dos etapas como la generación sin diagnóstico usan la misma instrucción
(Sección~\ref{sec:met_generacion}):
\begin{minted}[linenos=false]{text}
Describe this skin lesion.
\end{minted}

En la cascada y en el nivel afirmativo del ruteo, el generador recibe el diagnóstico con la plantilla
\emph{noreasoning} (Secciones~\ref{sec:met_cascada} y~\ref{sec:met_ruteo}):
\begin{minted}[linenos=false]{text}
The confirmed diagnosis is {dx}. Output ONLY the final clinical report — do NOT
include any reasoning, thinking steps, or <think> blocks. The report must
explicitly state the diagnosis ({dx}) and describe the supporting morphological
and dermatoscopic features. Do not mention alternative or differential
diagnoses.
\end{minted}

Las otras cuatro variantes comparadas en el piloto con que se eligió la anterior son las siguientes:
\begin{minted}[linenos=false]{text}
soft:
    The most likely diagnosis for this skin lesion is {dx}. Write a structured
    clinical report describing the lesion's morphology, dermatoscopic features
    and findings consistent with this diagnosis.
dxfirst:
    Diagnosis: {dx}. Write a clinical report that begins by stating this
    diagnosis ({dx}), then describes the dermatoscopic and morphological
    features that support it.
conclude:
    Write a structured clinical report for this skin lesion. You must end the
    report with this exact line: 'Most likely diagnosis: {dx}.'
confirmed:
    The confirmed diagnosis is {dx}. Write a concise clinical report that
    explicitly names this diagnosis ({dx}) and describes the supporting
    morphological and dermatoscopic features. Do not list alternative or
    differential diagnoses.
\end{minted}

En el nivel diferencial del ruteo, la plantilla presenta los dos candidatos, el diagnóstico del clasificador y el
voto del kNN, aunque atribuye la incertidumbre al clasificador:
\begin{minted}[linenos=false]{text}
The classifier is uncertain about this lesion; its most likely diagnoses are:
{candidates}. Output ONLY the final clinical report — no reasoning or <think>
blocks. Describe the morphological and dermatoscopic features, and present these
as a differential diagnosis (discuss the candidates) rather than committing to a
single one.
\end{minted}

Por su parte, las dos plantillas de recuperación (Sección~\ref{sec:met_recuperacion}) reciben como contexto una frase
con el diagnóstico de los tres vecinos y, entre paréntesis, su localización cuando el metadato la trae, por
ejemplo:
\begin{minted}[linenos=false]{text}
Retrieved visually similar dermatoscopic cases were diagnosed as: melanoma
(face); melanoma (face); actinic keratosis (face).
\end{minted}
La plantilla de recuperación sola pide el diagnóstico al generador, mientras que la de Swin más recuperación
le entrega el diagnóstico del clasificador:
\begin{minted}[linenos=false]{text}
rag:
    {rag_context} Based on this dermatoscopic image and the retrieved similar
    cases, write a clinical report describing the lesion's morphology and
    dermatoscopic features, and state the single most likely diagnosis.
grounded_rag:
    The confirmed diagnosis is {dx}. {rag_context} Output ONLY the final
    clinical report — no reasoning or <think> blocks. The report must explicitly
    state the diagnosis ({dx}) and describe the supporting morphological and
    dermatoscopic features. Do not list differential diagnoses.
\end{minted}

\section{Términos de las métricas de Nivel 2}

La consistencia diagnóstica reconoce cada clase con los términos de la Tabla~\ref{tab:ap_terminos_dx}
(Sección~\ref{sec:met_eval_generacion}). Todos los términos se buscan sin distinguir mayúsculas, y las abreviaturas y «mole» exigen palabra completa.

\begin{table}[H]
    \centering
    \caption{Términos con que la búsqueda estricta reconoce el diagnóstico de cada clase.}
    \label{tab:ap_terminos_dx}
    \small
    \begin{tabularx}{\textwidth}{@{}>{\raggedright\arraybackslash}p{2.2cm}>{\raggedright\arraybackslash}X@{}}
        \toprule
        \textbf{Clase} & \textbf{Términos} \\
        \midrule
        MEL & melanoma \\
        \filasep
        NV & nevus; nevi; mole (palabra completa) \\
        \filasep
        BCC & basal cell; basalioma; BCC (palabra completa) \\
        \filasep
        AKIEC & actinic kerato…; solar kerato…; AKIEC (palabra completa) \\
        \filasep
        BKL & benign kerato…; seborrh… kerato…; BKL (palabra completa) \\
        \filasep
        DF & dermatofibroma \\
        \filasep
        VASC & vascular lesion; hemangioma; angioma; pyogenic granuloma \\
        \bottomrule
    \end{tabularx}
\end{table}

La cobertura de entidades clínicas, en cambio, usa las tres categorías de la Tabla~\ref{tab:ap_entidades}, y cada
término se busca como palabra completa (Ecuación~\ref{eq:marco_cobertura_ent}).

\begin{table}[H]
    \centering
    \caption{Categorías y términos de la cobertura de entidades clínicas.}
    \label{tab:ap_entidades}
    \small
    \begin{tabularx}{\textwidth}{@{}>{\raggedright\arraybackslash}p{3.2cm}>{\raggedright\arraybackslash}X@{}}
        \toprule
        \textbf{Categoría} & \textbf{Términos} \\
        \midrule
        Criterios ABCDE & asymmetry; border; color; diameter; evolution \\
        \filasep
        Tipos de lesión & melanoma; nevus; basal cell carcinoma; actinic keratosis; benign keratosis;
            dermatofibroma; vascular lesion; papule; nodule; macule; plaque; patch \\
        \filasep
        Rasgos clínicos & pigmentation; ulceration; scaling; erythema; telangiectasia; regression; dermoscopic \\
        \bottomrule
    \end{tabularx}
\end{table}
